\section{Sharpness of the complete scalar family}

We now show that the constants above cannot be improved using only Proposition~\ref{prop:master} for all $c>0$ and the total-count identity.  This is a sharpness statement for the scalar relaxation, not a claim that the extremal distributions occur among actual zeta zeros.

Normalize total multiplicity to one.  Let $x_m\geq0$ denote the density of on-line locations of multiplicity $m$, and $y_m\geq0$ the density of off-line pairs of multiplicity $m$.  The abstract feasibility conditions are
\begin{equation}\label{eq:feasible}
\sum_{m\geq1}m x_m+2\sum_{m\geq1}m y_m=1
\end{equation}
and
\begin{equation}\label{eq:feasiblec}
2c-\kMT\leq\sum_{m\geq1}k_c(m)x_m+c^2\sum_{m\geq1}y_m
\qquad(c>0).
\end{equation}

\begin{theorem}[Sharpness]\label{thm:sharp}
For each fixed $r\geq1$, the lower bounds in Theorems~\ref{thm:multiset} and \ref{thm:distinct} are optimal over all nonnegative distributions satisfying \eqref{eq:feasible}--\eqref{eq:feasiblec}.
\end{theorem}

\begin{proof}
For $r=1$, take
\[
x_1=2-\kMT,
\qquad
x_2=\frac{\kMT-1}{2},
\qquad
x_m=y_m=0\ \text{otherwise}.
\]
Then \eqref{eq:feasible} holds.  If $0<c\leq1$, the slack in \eqref{eq:feasiblec} is decreasing and is nonnegative at $c=1$; for $1\leq c\leq2$ it equals
\[
\frac{\kMT-1}{2}(c-2)^2;
\]
for $c\geq2$ it is identically zero.  The low multiset mass is $2-\kMT$, and the low fraction among distinct locations is
\[
\frac{x_1}{x_1+x_2}=1-\frac{\kMT-1}{3-\kMT}.
\]

For $r=2$, take
\[
x_1=\frac{3-\kMT}{2},
\qquad
x_3=\frac{\kMT-1}{6},
\qquad
x_m=y_m=0\ \text{otherwise}.
\]
For $1\leq c\leq3$, the slack is
\[
\frac{\kMT-1}{6}(c-3)^2,
\]
while the ranges $c\leq1$ and $c\geq3$ are handled as above.  This realizes both $r=2$ bounds.

Now suppose $r\geq3$.  Put
\[
d:=\kMT-1,
\qquad
A:=3-2\sqrt2,
\]
and define
\begin{equation}\label{eq:sharpabc}
a:=1-d(\sqrt2-1),
\qquad
h:=\frac{Ad}{r-1},
\qquad
b:=\frac{1-a-(r+1)h}{2}.
\end{equation}
\Needspace{8\baselineskip}
All three numbers are nonnegative.  Indeed, $a>0$ because $0<d<1$ and $0<\sqrt2-1<1$, while $h>0$ and
\[
2b=d\left((\sqrt2-1)-(3-2\sqrt2)\frac{r+1}{r-1}\right)
\geq d(5\sqrt2-7)>0,
\]
since $(r+1)/(r-1)\leq2$ and $\sqrt2>7/5$.  Take
\[
x_1=a,
\qquad
x_{r+1}=h,
\qquad
y_1=b,
\]
with all other variables zero.  By construction,
\[
a+(r+1)h+2b=1.
\]
Let
\[
S(c):=a k_c(1)+h k_c(r+1)+bc^2-(2c-\kMT).
\]
For $1\leq c\leq r+1$, direct expansion using \eqref{eq:sharpabc} gives the exact square
\begin{equation}\label{eq:square}
S(c)=(h+b)(c-(2+\sqrt2))^2\geq0.
\end{equation}
For $0<c\leq1$,
\[
S(c)=(a+h+b)c^2-2c+\kMT.
\]
Since $a+h+b\leq1$, this is decreasing on $(0,1]$ and is nonnegative at $c=1$ by \eqref{eq:square}.  For $c\geq r+1$, differentiation and the total-count identity give
\[
S'(c)=2b(c-2)\geq0,
\]
so this range is minimized at $c=r+1$, where \eqref{eq:square} again applies.  Hence the distribution satisfies every constraint \eqref{eq:feasiblec}.

Its low multiset mass is
\[
a+2b
=1-Ad\frac{r+1}{r-1},
\]
which is exactly \eqref{eq:Mr}.  Its distinct total and high-multiplicity distinct mass are
\[
D=a+h+2b=1-\frac{Adr}{r-1},
\qquad
H_r=h=\frac{Ad}{r-1},
\]
so
\[
\frac{H_r}{D}=\frac{Ad}{r-1-rAd},
\]
which is equality in \eqref{eq:Dr}.  This proves sharpness.
\end{proof}

\begin{remark}[Why $2+\sqrt2$ appears]
For $r\geq3$, a simple on-line point has scalar charge $k_c(1)=2c-1$ when $c\geq1$, while each point in a simple off-line pair carries half of the pair charge, $c^2/2$.  Balancing those two worst cases gives
\[
2c-1=\frac{c^2}{2},
\]
whose relevant solution is $c=2+\sqrt2$.  The irrational parameter is therefore structural rather than a numerical accident.
\end{remark}
